\section{Conclusions}
\label{p:conclusions}

We have walked the discrete-to-continuous ladder
$\text{SSH} \to \text{graphene} \to \text{honeycomb QW} \to \text{Weyl/Dirac 3+1}$
keeping a uniform notation, three clearly separated levels (microscopic
exact, effective infrared, covariant spinorial), and a strict
verification policy: every algebraic and combinatorial claim is paired
with a symbolic test, and the build of the manuscript fails unless all
$146$ tests pass.

The relativistic structure --- the Clifford algebra of $\gamma^{\mu}$,
the spinorial Lorentz generators $\Sigma^{\mu\nu}$, the double cover by
$SL(2,\bbC)$, an effective Lorentzian light cone with controlled group
velocity, and a chirality balance that holds even in the absence of
translation invariance --- emerges entirely as an effective infrared
organization, not as a microscopic symmetry of the discrete step.

The two open problems identified in the long form of this work, on the
global spectral classification of the discrete step and on the
formulation of the underlying protospace without reciprocal-space
language, both come equipped with a first body of machine-checked
machinery, leaving them open in a precise rather than diffuse sense.

The verification methodology is, we believe, the most reusable
contribution. It does not require interactive theorem provers, integrates
with the scientist's existing toolchain, and yields manuscripts whose
algebraic claims are auditable in under a minute per claim.

\section*{Acknowledgements}
The author thanks the open-source scientific Python and TeX communities
whose tools made the verification and typesetting workflow possible.
AI-assisted coding and editorial tools were used to help draft, refactor,
and review parts of the validation modules, tests, and manuscript text;
all scientific claims, validator design choices, and final interpretations
remain the responsibility of the author.

\section*{Code and reproducibility}
The validation code is maintained in the companion project repository
under \texttt{validators/} and \texttt{tests/}. A public archival URL or
DOI should be inserted here before journal submission. The build pipeline
runs in under twenty seconds end-to-end on a standard laptop.
